%!TeX root=../princesstop.tex
\chapter{In the Attic}

\begin{figure}[t!]
\centering
\includegraphics[width=\linewidth]{filigree}
\end{figure}

\lettrine[lines=5]{T}{he} first night she spent in her attic was a thing Sara never forgot. During its passing she lived through a wild, unchildlike woe of which she never spoke to anyone about her. There was no one who would have understood. It was, indeed, well for her that as she lay awake in the darkness her mind was forcibly distracted, now and then, by the strangeness of her surroundings. It was, perhaps, well for her that she was reminded by her small body of material things. If this had not been so, the anguish of her young mind might have been too great for a child to bear. But, really, while the night was passing she scarcely knew that she had a body at all or remembered any other thing than one.

<My papa is dead!> she kept whispering to herself. <My papa is dead!>

It was not until long afterward that she realized that her bed had been so hard that she turned over and over in it to find a place to rest, that the darkness seemed more intense than any she had ever known, and that the wind howled over the roof among the chimneys like something which wailed aloud. Then there was something worse. This was certain scufflings and scratchings and squeakings in the walls and behind the skirting boards. She knew what they meant, because Becky had described them. They meant rats and mice who were either fighting with each other or playing together. Once or twice she even heard sharp-toed feet scurrying across the floor, and she remembered in those after days, when she recalled things, that when first she heard them she started up in bed and sat trembling, and when she lay down again covered her head with the bedclothes.

The change in her life did not come about gradually, but was made all at once.

<She must begin as she is to go on,> Miss Minchin said to Miss Amelia. <She must be taught at once what she is to expect.>

Mariette had left the house the next morning. The glimpse Sara caught of her sitting room, as she passed its open door, showed her that everything had been changed. Her ornaments and luxuries had been removed, and a bed had been placed in a corner to transform it into a new pupil's bedroom.

When she went down to breakfast she saw that her seat at Miss Minchin's side was occupied by Lavinia, and Miss Minchin spoke to her coldly.

<You will begin your new duties, Sara,> she said, <by taking your seat with the younger children at a smaller table. You must keep them quiet, and see that they behave well and do not waste their food. You ought to have been down earlier. Lottie has already upset her tea.>

That was the beginning, and from day to day the duties given to her were added to. She taught the younger children French and heard their other lessons, and these were the least of her labours. It was found that she could be made use of in numberless directions. She could be sent on errands at any time and in all weathers. She could be told to do things other people neglected. The cook and the housemaids took their tone from Miss Minchin, and rather enjoyed ordering about the <young one> who had been made so much fuss over for so long. They were not servants of the best class, and had neither good manners nor good tempers, and it was frequently convenient to have at hand someone on whom blame could be laid.

During the first month or two, Sara thought that her willingness to do things as well as she could, and her silence under reproof, might soften those who drove her so hard. In her proud little heart she wanted them to see that she was trying to earn her living and not accepting charity. But the time came when she saw that no one was softened at all; and the more willing she was to do as she was told, the more domineering and exacting careless housemaids became, and the more ready a scolding cook was to blame her.

If she had been older, Miss Minchin would have given her the bigger girls to teach and saved money by dismissing an instructress; but while she remained and looked like a child, she could be made more useful as a sort of little superior errand girl and maid of all work. An ordinary errand boy would not have been so clever and reliable. Sara could be trusted with difficult commissions and complicated messages. She could even go and pay bills, and she combined with this the ability to dust a room well and to set things in order.

Her own lessons became things of the past. She was taught nothing, and only after long and busy days spent in running here and there at everybody's orders was she grudgingly allowed to go into the deserted schoolroom, with a pile of old books, and study alone at night.

<If I do not remind myself of the things I have learned, perhaps I may forget them,> she said to herself. <I am almost a scullery maid, and if I am a scullery maid who knows nothing, I shall be like poor Becky. I wonder if I could \textsc{quite} forget and begin to drop my H's and not remember that Henry the Eighth had six wives.>

One of the most curious things in her new existence was her changed position among the pupils. Instead of being a sort of small royal personage among them, she no longer seemed to be one of their number at all. She was kept so constantly at work that she scarcely ever had an opportunity of speaking to any of them, and she could not avoid seeing that Miss Minchin preferred that she should live a life apart from that of the occupants of the schoolroom.

<I will not have her forming intimacies and talking to the other children,> that lady said. <Girls like a grievance, and if she begins to tell romantic stories about herself, she will become an ill-used heroine, and parents will be given a wrong impression. It is better that she should live a separate life—one suited to her circumstances. I am giving her a home, and that is more than she has any right to expect from me.>

Sara did not expect much, and was far too proud to try to continue to be intimate with girls who evidently felt rather awkward and uncertain about her. The fact was that Miss Minchin's pupils were a set of dull, matter-of-fact young people. They were accustomed to being rich and comfortable, and as Sara's frocks grew shorter and shabbier and queerer-looking, and it became an established fact that she wore shoes with holes in them and was sent out to buy groceries and carry them through the streets in a basket on her arm when the cook wanted them in a hurry, they felt rather as if, when they spoke to her, they were addressing an under servant.

<To think that she was the girl with the diamond mines,> Lavinia commented. <She does look an object. And she's queerer than ever. I never liked her much, but I can't bear that way she has now of looking at people without speaking—just as if she was finding them out.>

<I am,> said Sara, promptly, when she heard of this. <That's what I look at some people for. I like to know about them. I think them over afterward.>

The truth was that she had saved herself annoyance several times by keeping her eye on Lavinia, who was quite ready to make mischief, and would have been rather pleased to have made it for the ex-show pupil.

Sara never made any mischief herself, or interfered with anyone. She worked like a drudge; she tramped through the wet streets, carrying parcels and baskets; she laboured with the childish inattention of the little ones' French lessons; as she became shabbier and more forlorn-looking, she was told that she had better take her meals downstairs; she was treated as if she was nobody's concern, and her heart grew proud and sore, but she never told anyone what she felt.

<Soldiers don't complain,> she would say between her small, shut teeth, <I am not going to do it; I will pretend this is part of a war.>

But there were hours when her child heart might almost have broken with loneliness but for three people.

The first, it must be owned, was Becky—just Becky. Throughout all that first night spent in the garret, she had felt a vague comfort in knowing that on the other side of the wall in which the rats scuffled and squeaked there was another young human creature. And during the nights that followed the sense of comfort grew. They had little chance to speak to each other during the day. Each had her own tasks to perform, and any attempt at conversation would have been regarded as a tendency to loiter and lose time. <Don't mind me, miss,> Becky whispered during the first morning, <if I don't say nothin' polite. Some un'd be down on us if I did. I \textsc{means} <please> an' <thank you> an' <beg pardon,> but I dassn't to take time to say it.>

But before daybreak she used to slip into Sara's attic and button her dress and give her such help as she required before she went downstairs to light the kitchen fire. And when night came Sara always heard the humble knock at her door which meant that her handmaid was ready to help her again if she was needed. During the first weeks of her grief Sara felt as if she were too stupefied to talk, so it happened that some time passed before they saw each other much or exchanged visits. Becky's heart told her that it was best that people in trouble should be left alone.

The second of the trio of comforters was Ermengarde, but odd things happened before Ermengarde found her place.

When Sara's mind seemed to awaken again to the life about her, she realized that she had forgotten that an Ermengarde lived in the world. The two had always been friends, but Sara had felt as if she were years the older. It could not be contested that Ermengarde was as dull as she was affectionate. She clung to Sara in a simple, helpless way; she brought her lessons to her that she might be helped; she listened to her every word and besieged her with requests for stories. But she had nothing interesting to say herself, and she loathed books of every description. She was, in fact, not a person one would remember when one was caught in the storm of a great trouble, and Sara forgot her.

It had been all the easier to forget her because she had been suddenly called home for a few weeks. When she came back she did not see Sara for a day or two, and when she met her for the first time she encountered her coming down a corridor with her arms full of garments which were to be taken downstairs to be mended. Sara herself had already been taught to mend them. She looked pale and unlike herself, and she was attired in the queer, outgrown frock whose shortness showed so much thin black leg.

Ermengarde was too slow a girl to be equal to such a situation. She could not think of anything to say. She knew what had happened, but, somehow, she had never imagined Sara could look like this—so odd and poor and almost like a servant. It made her quite miserable, and she could do nothing but break into a short hysterical laugh and exclaim—aimlessly and as if without any meaning, <Oh, Sara, is that you?>

<Yes,> answered Sara, and suddenly a strange thought passed through her mind and made her face flush. She held the pile of garments in her arms, and her chin rested upon the top of it to keep it steady. Something in the look of her straight-gazing eyes made Ermengarde lose her wits still more. She felt as if Sara had changed into a new kind of girl, and she had never known her before. Perhaps it was because she had suddenly grown poor and had to mend things and work like Becky.

<Oh,> she stammered. <How—how are you?>

<I don't know,> Sara replied. <How are you?>

<I'm—I'm quite well,> said Ermengarde, overwhelmed with shyness. Then spasmodically she thought of something to say which seemed more intimate. <Are you—are you very unhappy?> she said in a rush.

Then Sara was guilty of an injustice. Just at that moment her torn heart swelled within her, and she felt that if anyone was as stupid as that, one had better get away from her.

<What do you think?> she said. <Do you think I am very happy?> And she marched past her without another word.

In course of time she realized that if her wretchedness had not made her forget things, she would have known that poor, dull Ermengarde was not to be blamed for her unready, awkward ways. She was always awkward, and the more she felt, the more stupid she was given to being.

But the sudden thought which had flashed upon her had made her over-sensitive.

<She is like the others,> she had thought. <She does not really want to talk to me. She knows no one does.>

So for several weeks a barrier stood between them. When they met by chance Sara looked the other way, and Ermengarde felt too stiff and embarrassed to speak. Sometimes they nodded to each other in passing, but there were times when they did not even exchange a greeting.

<If she would rather not talk to me,> Sara thought, <I will keep out of her way. Miss Minchin makes that easy enough.>

Miss Minchin made it so easy that at last they scarcely saw each other at all. At that time it was noticed that Ermengarde was more stupid than ever, and that she looked listless and unhappy. She used to sit in the window-seat, huddled in a heap, and stare out of the window without speaking. Once Jessie, who was passing, stopped to look at her curiously.

<What are you crying for, Ermengarde?> she asked.

<I'm not crying,> answered Ermengarde, in a muffled, unsteady voice.

<You are,> said Jessie. <A great big tear just rolled down the bridge of your nose and dropped off at the end of it. And there goes another.>

<Well,> said Ermengarde, <I'm miserable—and no one need interfere.> And she turned her plump back and took out her handkerchief and boldly hid her face in it.

That night, when Sara went to her attic, she was later than usual. She had been kept at work until after the hour at which the pupils went to bed, and after that she had gone to her lessons in the lonely schoolroom. When she reached the top of the stairs, she was surprised to see a glimmer of light coming from under the attic door.

<Nobody goes there but myself,> she thought quickly, <but someone has lighted a candle.>

Someone had, indeed, lighted a candle, and it was not burning in the kitchen candlestick she was expected to use, but in one of those belonging to the pupils' bedrooms. The someone was sitting upon the battered footstool, and was dressed in her nightgown and wrapped up in a red shawl. It was Ermengarde.

<Ermengarde!> cried Sara. She was so startled that she was almost frightened. <You will get into trouble.>

Ermengarde stumbled up from her footstool. She shuffled across the attic in her bedroom slippers, which were too large for her. Her eyes and nose were pink with crying.

<I know I shall—if I'm found out.> she said. <But I don't care—I don't care a bit. Oh, Sara, please tell me. What is the matter? Why don't you like me any more?>

Something in her voice made the familiar lump rise in Sara's throat. It was so affectionate and simple—so like the old Ermengarde who had asked her to be <best friends.> It sounded as if she had not meant what she had seemed to mean during these past weeks.

<I do like you,> Sara answered. <I thought—you see, everything is different now. I thought you—were different.>

Ermengarde opened her wet eyes wide.

<Why, it was you who were different!> she cried. <You didn't want to talk to me. I didn't know what to do. It was you who were different after I came back.>

Sara thought a moment. She saw she had made a mistake.

<I \textsc{am} different,> she explained, <though not in the way you think. Miss Minchin does not want me to talk to the girls. Most of them don't want to talk to me. I thought—perhaps—you didn't. So I tried to keep out of your way.>

<Oh, Sara,> Ermengarde almost wailed in her reproachful dismay. And then after one more look they rushed into each other's arms. It must be confessed that Sara's small black head lay for some minutes on the shoulder covered by the red shawl. When Ermengarde had seemed to desert her, she had felt horribly lonely.

Afterward they sat down upon the floor together, Sara clasping her knees with her arms, and Ermengarde rolled up in her shawl. Ermengarde looked at the odd, big-eyed little face adoringly.

<I couldn't bear it any more,> she said. <I dare say you could live without me, Sara; but I couldn't live without you. I was nearly \textsc{dead}. So tonight, when I was crying under the bedclothes, I thought all at once of creeping up here and just begging you to let us be friends again.>

<You are nicer than I am,> said Sara. <I was too proud to try and make friends. You see, now that trials have come, they have shown that I am \textsc{not} a nice child. I was afraid they would. Perhaps>—wrinkling her forehead wisely—<that is what they were sent for.>

<I don't see any good in them,> said Ermengarde stoutly.

<Neither do I—to speak the truth,> admitted Sara, frankly. <But I suppose there \textsc{might} be good in things, even if we don't see it. There \textsc{might}>—doubtfully—<be good in Miss Minchin.>

Ermengarde looked round the attic with a rather fearsome curiosity.

<Sara,> she said, <do you think you can bear living here?>

Sara looked round also.

<If I pretend it's quite different, I can,> she answered; <or if I pretend it is a place in a story.>

She spoke slowly. Her imagination was beginning to work for her. It had not worked for her at all since her troubles had come upon her. She had felt as if it had been stunned.

<Other people have lived in worse places. Think of the Count of Monte Cristo in the dungeons of the Château d'If. And think of the people in the Bastille!>

<The Bastille,> half whispered Ermengarde, watching her and beginning to be fascinated. She remembered stories of the French Revolution which Sara had been able to fix in her mind by her dramatic relation of them. No one but Sara could have done it.

A well-known glow came into Sara's eyes.

<Yes,> she said, hugging her knees, <that will be a good place to pretend about. I am a prisoner in the Bastille. I have been here for years and years—and years; and everybody has forgotten about me. Miss Minchin is the jailer—and Becky>—a sudden light adding itself to the glow in her eyes—<Becky is the prisoner in the next cell.>

She turned to Ermengarde, looking quite like the old Sara.

<I shall pretend that,> she said; <and it will be a great comfort.>

Ermengarde was at once enraptured and awed.

<And will you tell me all about it?> she said. <May I creep up here at night, whenever it is safe, and hear the things you have made up in the day? It will seem as if we were more <best friends> than ever.>

<Yes,> answered Sara, nodding. <Adversity tries people, and mine has tried you and proved how nice you are.>